\itemhead{复习题一 · 第1题（教材p.29）}{XJ-B1-C01-RV-P029-Y0406}
\label{XJ-B1-C01-RV-P029-Y0406}
{\footnotesize 湘教版教材p.29 / PDF第36页 / 复习题一 / 第1题}
\field{原题}{下列集合表示同一集合的是（\quad）
\choice{A}$M=\{1,2\}$，$N=\{(1,2)\}$
\choice{B}$M=\{(2,1)\}$，$N=\{(1,2)\}$
\choice{C}$M=\{2,1\}$，$N=\{1,2\}$
\choice{D}$M=\{x\mid x+2y=3\}$，$N=\{y\mid x+2y=3,\ x\in\mathbb Z\}$}
\field{参考答案}{C。集合元素的次序无关；数与有序数对不能混同。}
{\small \field{关联知识}{S02 集合表示、区间与数集点集（教材帮p.2、3、4 / PDF11、12、13）}
\field{方法与易错点}{先求出约束的解，再选列举法或描述法，必要时在数轴或坐标系中表示。 易错：遗漏0或负整数；将一个有序数对拆成两个元素；区间端点误取。}
\field{答案依据}{答案册相关题组：PDF第8、9页（册内页7、8）；按当前题库原题逐项重算，不按旧版题号机械对套。}}
